\documentclass[aspectratio=169]{beamer}
\usepackage{spc}
\usepackage{textpos}
\begin{document}

\begin{frame}
  \title{\vspace{-5ex}\darkblue Yellowfin Model Development\\[2ex]
    \it\large\darkgray First Runs and Ideas to Explore}
  \author{\vspace{-10ex}\darkgray\bf Arni Magnusson}
  \date{\darkgreen SPC PAW Meeting\\[0.5ex]
    25 March 2026, Nouméa}
  \titlepage
  \begin{textblock}{0}(10,0.5)
    \includegraphics[width=4cm]{fish_yellowfin}
  \end{textblock}
\end{frame}

% ______________________________________________________________________________

\begin{frame}{Model development}
  Alternative CPUE indices
  Time-varying CPUE index CV
  Richards Growth
  Convert weight comps to length comps
  Combine fisheries
  Lengh-based selectivity
  Growth, three levels of k
  Natural mortality, three levels
  SEAPODYM movement coefficients
  Tag mixing for each release group
  Revisit composition weights: possibly downweighting temperate size comps
  Revisit composition weights: lower likelihood for size comps in general
  Bigeye regional structure
  Single-region or equatorial model?
  2023 PAW slide
  • Analyze data conflict, identify and test different approaches
  • Indonesian catches
  • Alternative CPUE indices
  • Release group-specific tag mixing rate
\end{frame}

% ______________________________________________________________________________

\begin{frame}
  \vspace{4ex}
  These slides:\\[2ex]
  \begin{center}
    \blue\url{https://github.com/PacificCommunity/ofp-sam-yft-2026-assessment}
  \end{center}
\end{frame}

\end{document}
